\setcounter{page}{1}
\pagestyle{body}
\hspace*{\fill}
\section*{\zihao{-2}\heiti
基于A*算法的智能贪吃蛇游戏设计与实现}

\noindent
{\zihao{-4}\heiti
摘要：}
贪吃蛇游戏自诞生以来便是一款经久不衰的经典小游戏，其规则简单但策略空间丰富，适合用来验证各类搜索算法的实际效果。本文基于Python语言和Pygame图形库，设计并实现了一款具备自动寻路能力的智能贪吃蛇游戏。系统的核心在于利用A*算法驱动蛇的路径规划，使蛇能够在避开自身蛇身与障碍物的前提下，持续追踪食物位置并做出移动决策。论文从A*算法的基本原理出发，分析了曼哈顿距离、欧几里得距离等启发式函数对搜索效率和路径质量的影响，并将其与BFS、Dijkstra等经典搜索算法进行了性能对比。在实际调试过程中发现，启发式函数的权重选择对搜索结果有显著影响——权重过大会导致搜索不够彻底，错过最优路径；权重过小则退化为类似Dijkstra的遍历，计算开销陡增。此外，当蛇身变长导致可用空间急剧缩减时，单纯的A*寻路往往无法避免"自己堵死自己"的死局，需要引入安全区域评估等辅助策略。本系统实现了手动操控与AI自动运行两种模式，支持游戏界面可视化、分数统计、速度调节等功能。测试结果表明，在中等规模网格下A*算法能够以较低的搜索代价完成绝大多数场景的路径规划，但在高密度蛇身条件下仍存在决策失误的情况，这也是后续优化的方向。

\noindent{\heiti 关键词：}A*算法；贪吃蛇游戏；路径规划；Pygame；游戏AI

\clearpage
\hspace*{\fill}
\section*{\zihao{-2}
Design and Implementation of an Intelligent Snake Game Based on A* Algorithm}

\noindent\textbf{Abstract}: 
The snake game has been a classic arcade game since its inception, featuring simple rules yet offering rich strategic depth, making it an excellent testbed for various search algorithms. This thesis designs and implements an intelligent snake game with automatic pathfinding capability, built upon the Python programming language and the Pygame graphical library. The core of the system lies in employing the A* algorithm to drive path planning for the snake, enabling it to continuously track food positions and make movement decisions while avoiding collisions with its own body and obstacles. Starting from the fundamental principles of the A* algorithm, the thesis analyzes the impact of heuristic functions such as Manhattan distance and Euclidean distance on search efficiency and path quality, and compares their performance against classical search algorithms including BFS and Dijkstra. During practical debugging, it was found that the weight selection of heuristic functions significantly affects search results --- excessively large weights lead to insufficient search and suboptimal paths, while overly small weights degenerate into Dijkstra-like traversal with sharply increased computational overhead. Furthermore, when the snake body grows and available space diminishes, pure A* pathfinding often fails to prevent self-blockage, necessitating auxiliary strategies such as safe area evaluation. The system implements both manual control and AI automatic operation modes, supporting game interface visualization, score statistics, and speed adjustment. Test results demonstrate that the A* algorithm can accomplish path planning in most scenarios with low search cost under moderate grid sizes, yet decision errors still occur under high-density snake body conditions, pointing toward directions for future optimization.

\noindent\textbf{Keywords：}A* Algorithm; Snake Game; Path Planning; Pygame; Game AI
